%%%%%%%%%%%%%%%%%%%%%%%%%%%%%%%%%%%%%%%%%%%%%%%%%%%%%%%%%%%%%%%%%%%%%%%%%%%%%%%%%%%
\subsection{Trigger Efficiencies~\label{sec:medium_muon_trigEff}}
\begin{figure}[h] 
\centering
\includegraphics[width=0.5\textwidth]%
{figures/AllFigs/Can_trigeff_type1_trig0_MuonSelectionCut0}%
%\includegraphics[width=0.5\textwidth]%
%{figures/AllFigs/Can_trigeff_type1_trig0_MuonSelectionCut1}
\caption{
HLT\_mu4 trigger efficiency for \PbPb.
Plots correspond to Medium muons.
\label{fig:pbpb_trigEff2}
}
\end{figure}


\begin{figure}[h]
\centering
\includegraphics[width=1\textwidth]%
{figures/AllFigs/Can_eff_trig_option0_MuonSelectionCut0}%
\caption{
HLT\_mu4 trigger efficiency for \PbPb.
The lines indicate fits to Fermi+linear functions given in Eq.~\eqref{eq:trigger_eff_fit}.
The fits are made over the range \pt>4~\GeV.
The plots correspond to Medium muons.
\label{fig:pbpb_trigEff_1D1}
}
\end{figure}


\begin{figure}
\centering
\includegraphics[width=0.5\textwidth]%
{figures/AllFigs/can_eff_Mu4NoL1_RefAll_MuonSelectionCut0}%
\includegraphics[width=0.5\textwidth]%
{figures/AllFigs/can_eff_Both_RefAll_MuonSelectionCut0}
\caption{
  Left: HLT\_mu4noL1 trigger efficiency for \PbPb.
  Right: HLT\_mu4noL1$\cap$HLT\_mu4 Triggers trigger efficiency for \PbPb.
  Plots correspond to Medium muons.
\label{fig:pbpb_trigEff_noL1_medium}
}
\end{figure}


\begin{figure}
\centering
\includegraphics[width=1\textwidth]%
{figures/AllFigs/can_eff1D_All_RefAll_MuonSelectionCut0}%
\caption{
  Efficiencies for the HLT\_mu4NoL1 trigger (red) as well as the combined 
    HLT\_mu4noL1$\cap$HLT\_mu4 eficiency (blue, labelled as \textit{Both}) in \PbPb.
  Each panel corresponds to a different $q\eta$ interval.
  The efficiency for the HLT\_mu4 trigger obtained using the HP
    stream is also shown (black).
  The plots correspond to Medium muons.
\label{fig:pbpb_trigEff_1D2_noL1_medium}
}
\end{figure}


\begin{figure}
\centering
\includegraphics[width=1\textwidth]%
{figures/AllFigs/can_eff1D_Mu4_RefAll_MuonSelectionCut1}%
\caption{
  Comparison of the efficiency for the HLT\_mu4 trigger obtained using the HP
    stream to that using the min-bias events.
  The plots correspond to Medium muons.
\label{fig:pbpb_trigEff_1D1_compare_medium}
}
\end{figure}


\FloatBarrier
\clearpage
\subsection{Reconstruction Efficienciess \label{sec:medium_muon_recoeff}}
\begin{figure}[h]
\centering
\includegraphics[width=0.5\textwidth]%
{figures/AllFigs/Can_recoeff2D_medium}%
%\includegraphics[width=0.5\textwidth]%
%{figures/AllFigs/Can_recoeff2D_tight}
\caption{
Muon reconstruction efficiencies for \PbPb.
The plot corresponds to Medium muons.
\label{fig:pbpb_ReconEff_2D}
}
\end{figure}


\begin{figure}[h]
\centering
\includegraphics[width=1.0\textwidth]%
{figures/AllFigs/Can_recoeff1D_medium}
\caption{
Muon reconstruction efficiency as a function of \pt and
  for several $q*\eta$ slices.
The lines indicate fits to Fermi functions.
Plots are for Medium quality muons.
\label{fig:pbpb_ReconEff}
}
\end{figure}

%%%%%%%%%%%%%%%%%%%%%%%%%%%%%%%%%%%%%%%%%%%%%%%%%%%%%%%%%%%%%%%%%%%%%%%%%%%%%%%%%%%

\FloatBarrier
\clearpage
\subsection{Fits to the \Dphi\ distributions}
\begin{figure}[h]
\centering
\includegraphics[width=1.0\textwidth]%
{figures/AllFigs/can_dphi_fits_cent12_LorentzZYAM_pTBar_GT5_EffCor}
\includegraphics[width=1.0\textwidth]%
{figures/AllFigs/can_dphi_fits_cent13_LorentzZYAM_pTBar_GT5_EffCor}
\includegraphics[width=1.0\textwidth]%
{figures/AllFigs/can_dphi_fits_cent4_LorentzZYAM_pTBar_GT5_EffCor}
\includegraphics[width=1.0\textwidth]%
{figures/AllFigs/can_dphi_fits_cent6_LorentzZYAM_pTBar_GT5_EffCor}
\includegraphics[width=1.0\textwidth]%
{figures/AllFigs/can_dphi_fits_cent8_LorentzZYAM_pTBar_GT5_EffCor}
\caption{
Fits to the \Dphi\ distributions using Eq.~\eqref{eq:fit_func}.
Plots are for the \PbPb\ data and from the Medium working point.
The left, middle and right panels correspond to same-charge, 
  opposite-charge and combined-charge pairs, respectively. 
From top to bottom, each row corresponds to a different centrality interval.
Note: the $y$-range of the different panels are very different,
  and they are zero suppressed.
\label{fig:fit_PbPb_medium}
}
\end{figure}



%%%%%%%%%%%%%%%%%%%%%%%%%%%%%%%%%%%%%%%%%%%%%%%%%%%%%%%%%%%%%%%%%%%%%%%%%%%%%%%%%%%
\begin{comment}
\clearpage
\subsection{Template fits for medium muons\label{sec:MediumMuons}}

Template-fit for Medium muon selection, \ptab>4~\GeV, and for
  several centrality intervals in Figure~\ref{fig:Templatefits_bycent_qual0_pt4p0to15p0_sign2}
Figure~\ref{fig:Signalfractions_qual0}, shows the \textit{signal-fractions} as a function
  of centrality for Medium muons. 

\begin{figure}[h]
\centering
\includegraphics[width=1.0\textwidth]%                                                                                           
{figures/DpopAna/Templatefits_bycent_qual0_pt4p0to15p0_sign2}
\caption{
Template fits to data (in black) for the $\Delta p/p$ pair significance. 
The   \textbf{Signal-Signal} contribution is shown in green, 
  the \textbf{Signal-Background} in magenta and 
  the \textbf{Background-Background} in blue.
The combined fit is shown in red.
The individual panels correspond to different centrality intervals.
The fits are for \ptab>4~\GeV, for Medium muons, and for the combined-charge pairs.
\label{fig:Templatefits_bycent_qual0_pt4p0to15p0_sign2}
}
\end{figure}



\begin{figure}[h]
\centering
\includegraphics[width=1.0\textwidth]%
{figures/DpopAna/Signalfractions_qual0}
\caption{
Signal fractions obtained from the template fits for same-sign muon pairs (open circles), 
  opposite-sign muon pairs (open squares), and their combination (closed circles) 
  as a function of centrality.
Plots are for Medium muons.
\label{fig:Signalfractions_qual0}
}
\end{figure}
\end{comment}
%%%%%%%%%%%%%%%%%%%%%%%%%%%%%%%%%%%%%%%%%%%%%%%%%%%%%%%%%%%%%%%%%%%%%%%%%%%%%%%%%%%


%%%%%%%%%%%%%%%%%%%%%%%%%%%%%%%%%%%%%%%%%%%%%%%%%%%%%%%%%%%%%%%%%%%%%%%%%%%%%%%%%%%
\clearpage
\subsection{Template fits for medium muons\label{sec:MediumMuons}}
This section discusses the results of the Template-fits to the $\dpop$ distributions
  for Medium muons.
Template-fit for several centrality intervals are shown in Figure~\ref{fig:Templatefits_ch2_dphi3_qual0}.
The fits are made for the difference between the Away-side ($|\Dphi|>\pi/2$) 
  and Same-size ($|\Dphi|<\pi/2$) pairs. 
The difference removes the contribution from the combinatorial pairs and thus gives
  the signal fraction only for the correlated pairs.
Figure~\ref{fig:Signalfractions_ch2_qual0}, shows the \textit{signal-fractions} as a function
  of centrality. 
The fractions are shown for pairs on the 
\begin{enumerate}
  \item Same-size ($|\Dphi|<\pi/2$), 
  \item Away-side ($|\Dphi|>\pi/2$),
  \item All-Sides (no restriction on $\Dphi$)
  \item The difference of the Away-side and Same-side distribution.
\end{enumerate}
The average signal fractions for the difference is $\sim0.85$ (obtained from the fit show in Figure~\ref{fig:Signalfractions_ch2_qual0}).
The signal factions are smaller for either the Same-sign or Oppodite-sign pairs, compared to the
  difference between the two, which indicates that the correlated pairs have a much larger
  signal fraction compared to the combinatorial contribution.



\begin{figure}[h]
\centering
\includegraphics[width=1.0\textwidth]%                                                                                           
{figures/AllFigs/can_fits_ch2_dphi3_templatefittype0_option0_qual0_pTBar_GT5_EffCor}
\caption{
  Template fits to the $\dpop$ pair significance. 
  The individual panels correspond to different centrality intervals.
  The fits are for the combined-charge pairs.
  The fits are made for the  difference between the Away-side ($|\Dphi|>\pi/2$) and Same-size ($|\Dphi|<\pi/2$) pairs. 
\label{fig:Templatefits_ch2_dphi3_qual0}
}
\end{figure}



\begin{figure}[h]
\centering
\includegraphics[width=1.0\textwidth]%
{figures/AllFigs/can_sig_frac_compare_dphi_ch2_templatefittype0_option0_qual0_pTBar_GT5_EffCor}
\caption{
  Signal fractions obtained from the template fits. 
  The fractions are shown for pairs on the 
    1) Same-size ($|\Dphi|<\pi/2$), 
    2) Away-side ($|\Dphi|>\pi/2$),
    3) All-Sides (no restriction on $\Dphi$), and
    4) the difference of the Away-side and Same-side distribution.
  The horizontal dashed line is a fit to the difference distribution.
  Plots are for Medium muons and All-charge pairs.
\label{fig:Signalfractions_ch2_qual0}
}
\end{figure}


\begin{comment}
\begin{figure}[h]
\centering
\includegraphics[width=1.0\textwidth]%
{figures/AllFigs/can_sig_frac_compare_dphi_ch0_templatefittype0_option0_qual0_pTBar_GT5_EffCor}
\caption{
Same as Figure~\ref{fig:Signalfractions_ch2_qual0} but for same-charge pairs.
\label{fig:Signalfractions_ch0_qual0}
}
\end{figure}

\begin{figure}[h]
\centering
\includegraphics[width=1.0\textwidth]%
{figures/AllFigs/can_sig_frac_compare_dphi_ch1_templatefittype0_option0_qual0_pTBar_GT5_EffCor}
\caption{
Same as Figure~\ref{fig:Signalfractions_ch2_qual0} but for opposite-charge pairs.
\label{fig:Signalfractions_ch1_qual0}
}
\end{figure}
\end{comment}
%%%%%%%%%%%%%%%%%%%%%%%%%%%%%%%%%%%%%%%%%%%%%%%%%%%%%%%%%%%%%%%%%%%%%%%%%%%%%%%%%%%
